\section{Introduction}
\label{sec:introduction}

Large language models (LLMs) are rapidly replacing human annotators in computational social science~\cite{gilardi2023chatgpt, ziems2024can, egami_hinck_stewart_wei2023, baumann2025}. Most tasks are inherently multi-class (stance detection: 3--4 categories, sentiment: 5, disaster response: 10+), and annotation errors propagate into downstream causal estimates in ways that depend on both error structure and causal design. \citet{baumann2025} show that 31\% of hypotheses yield incorrect conclusions with state-of-the-art LLMs, with 60\% error correlation across models~\cite{kim_garg2025}. Standard metrics do not capture this risk: within $F_1$ 0.61--0.69, causal bias varies by an order of magnitude across topologies.

The transition at $K{\geq}3$ is not gradual but structural. For $K{=}2$, the stochastic constraint forces any confusion matrix with positive diagonal to preserve coefficient signs, an algebraic property of the $2{\times}2$ matrix dimension~\cite{brenner1993}. This guarantee vanishes at $K{\geq}3$: the higher-dimensional confusion matrix admits negative effective weights on categories with opposite-sign coefficients; sign reversal can occur even under high accuracy~\cite{veirod_laake2001}. Under a uniform Dirichlet prior, 83.4\% of $K{=}3$ exposure confusion matrices produce sign reversal of at least one secondary coefficient; under realistic accuracy (diagonal $\geq 0.7$), 2.1\% still do. The primary treatment coefficient is never reversed under the tested $\beta$ configurations where a single category dominates ($\beta_1 \gg \beta_2$), an empirical regularity we discuss in Section~\ref{sec:method}.

$F_1$ does not exhibit significant rank correlation with causal bias ($\rho = -0.10$, $p = 0.45$ across 56 HumAID disaster-response scenarios ($K{=}10$); bias varies $>$10${\times}$ (from $<$6\% to 42.1\%) within $F_1$ 0.614--0.690). Prior work noted polytomous sign reversal~\cite{dosemeci1990, veirod_laake2001} and confounder amplification~\cite{brenner1993}, but covered at most one topology~\cite{lee_wood_doughty2024}. We present the first systematic analysis jointly addressing multi-class ($K{\geq}3$) misclassification across seven directed acyclic graph (DAG) topologies, with analytical predictions matching Monte Carlo within 1.6\% mean absolute percentage error (MAPE) across 161 scenarios.

We address these gaps with three contributions:

\begin{enumerate}
\item \textbf{Multi-class bias taxonomy.} We derive bias expressions for seven canonical DAG topologies\footnote{These span the four structural roles of the annotated variable: treatment (exposure), outcome-side covariate (confounding, mediation, M-bias, front-door), selection variable (collider), and instrumental variable (IV).} (Figure~\ref{fig:dag_taxonomy}) under $K$-class non-differential misclassification and establish a sign-preservation/reversal dichotomy: five topologies preserve the treatment effect sign under topology-specific regularity conditions (Proposition~\ref{prop:signflip}), while exposure and IV can reverse it.

\item \textbf{Annotation Sensitivity Values.} We define three ASV diagnostics (magnitude, significance, and sign-flip) with proved envelope monotonicity (Theorems~\ref{thm:monotone},~\ref{thm:k3_symmetric}) and $\varepsilon$-approximate monotonicity for general $K{\geq}3$ (Conjecture~\ref{conj:monotone}). This reveals a vulnerability spectrum from IV (maximally fragile) to M-bias (immune).

\item \textbf{Topology-aware correction routing.} Frisch--Waugh--Lovell (FWL) invariance renders matrix-based corrections provably ineffective for four topologies (confounding, mediation, collider, M-bias). We evaluate five correction methods and provide a topology-aware decision procedure (Algorithm~\ref{alg:correction}), validated with real LLM confusion matrices.
\end{enumerate}
